\documentclass{article}
\usepackage[T1]{fontenc}
\usepackage{lmodern}
\usepackage{graphicx}
\usepackage{amsmath,amssymb}
\usepackage[a4paper,margin=2cm]{geometry}
\usepackage[absolute,overlay]{textpos}
\setlength{\TPHorizModule}{1mm}
\setlength{\TPVertModule}{1mm}
\usepackage[indonesian]{babel}
\usepackage{hyperref}
\setlength{\emergencystretch}{2em}
% unit: Halaman 3

\begin{document}

% === Halaman 3 (python/native) ===

3

• Meskipun disebut kriptografi modern, namun algoritmanya tetap 

menggunakan dua teknik dasar dasar  di dalam kriptografi klasik: teknik 

substitusi dan teknik transposisi (permutasi), 

• tetapi operasinya dibuat lebih kompleks, tidak sesederhana cipher klasik.

   Tujuannya: agar cipher modern lebih sulit dikriptanalisis

• Selain kedua teknik dasar tersebut, juga digunakan teknik lain seperti 

rotasi, kompresi, ekspansi, penjumlahan modulo, dan lain-lain.

• Kriptografi modern melahirkan konsep-konsep baru seperti algoritma  

kriptografi kunci-publik, fungsi hash, protokol kriptografi, tanda-tangan 

digital, pembangkit bilangan acak, skema pembagian kunci, dsb.

\end{document}
